\section{The first toric layer}
\label{sec:toric}

Fix \(i\) and a field \(k\).  Let
\[
 S_i=k[x_B:B\in\cB_i]
\]
and define the incidence monomial map
\[
 \pi_i:S_i\longrightarrow k[y_p:p\in\cP],
 \qquad
 x_B\longmapsto\prod_{p\in\operatorname{PG}(B)}y_p.
\]
The incidence toric ideal is \(I_i=\ker\pi_i\).  Since every block
contains \(15\) points, \(I_i\) is homogeneous for total degree.

\begin{lemma}[No linear or quadratic moves]
\label{lem:no-small-toric}
The ideal \(I_i\) has no nonzero binomial of degree one or two.
\end{lemma}

\begin{proof}
Distinct blocks have distinct point-incidence columns, excluding degree
one.  Suppose a quadratic relation remains after cancelling common
variables, and choose a block \(B\) from one side.  If that side contains
\(B^2\), its incidence vector has scalar product at least \(30\) with
the column of \(B\).  Otherwise it contains \(B\) once and the scalar
product is at least \(15\).  A different four-space meets \(B\) in at
most \(7\) projective points, so the two blocks on the opposite side
contribute at most \(14\).  Both cases are impossible.
\end{proof}

\begin{lemma}[Squarefree cubic classification]
\label{lem:cubic-classification}
Every nonzero cubic binomial in \(I_i\), after cancellation, is the
binomial of a minimum opposite-regulus trade.
\end{lemma}

\begin{proof}
If a block \(B\) is repeated on one side, its scalar product with that
side is at least \(30\).  The three distinct blocks on the other side
contribute at most \(3\cdot7=21\).  Thus both cubic monomials are
squarefree.  A block shared by the two sides would cancel to a nonzero
quadratic, contradicting \cref{lem:no-small-toric}.  Hence the binomial
has three blocks in each leg and total support six.  It is a minimum
\(T_2(1,4,8)\) trade, and Krotov's \(t=1\) classification identifies it
with an opposite-regulus trade \cite{Krotov2017}.
\end{proof}

For a monomial \(M\), let
\[
 \mathcal F(M)=\{N:\pi_i(N)=\pi_i(M)\}
\]
be its multidegree fiber.

\begin{lemma}[Two-element fibers]
\label{lem:two-element-fiber}
If \(M=x_{B_1}x_{B_2}x_{B_3}\) is one leg of a minimum trade and
\(N\) is the other leg, then
\[
 \mathcal F(M)=\{M,N\}.
\]
\end{lemma}

\begin{proof}
Every monomial in the fiber has degree three.  By
\cref{lem:no-small-toric,lem:cubic-classification}, a second monomial is
squarefree, shares no block with \(M\), and forms a minimum regulus trade
with it.

The three blocks in \(M\) recover
\[
 Z=B_1\cap B_2\cap B_3,
 \qquad
 U=B_1+B_2+B_3.
\]
Their images in \(\operatorname{PG}(U/Z)\cong\operatorname{PG}(3,2)\)
are three skew lines.  Those lines have a unique opposite regulus, which
lifts to the three blocks of \(N\).  Thus no third monomial occurs.
\end{proof}

\begin{theorem}[Complete cubic layer]
\label{thm:toric-cubics}
For each \(i\), the primitive cubics of \(I_i\) are exactly the minimum
opposite-regulus binomials.  Every one is indispensable.  Their numbers
are
\[
 2\,271\,676,\qquad
 2\,213\,026,\qquad
 2\,216\,868
\]
for \(i=1,2,3\), respectively.  Equivalently,
\[
 \beta_{0,3}(I_i)=\beta_{1,3}(S_i/I_i)=|\cC_i|.
\]
\end{theorem}

\begin{proof}
\Cref{lem:cubic-classification} identifies all cubics, and
\cref{lem:two-element-fiber} shows that each has a two-element fiber.
The indispensable-fiber criterion of Aoki--Takemura--Yoshida implies that
the corresponding binomial occurs, up to scalar, in every Markov basis
\cite{AokiTakemuraYoshida}.  Since
\cref{lem:no-small-toric} excludes lower degrees, these cubics are minimal
generators.  The exact counts are those of
\cref{prop:circuit-census}.
\end{proof}

\begin{remark}[Higher degrees]
The theorem determines the first nonzero graded layer.  It does not prove
that \(I_i\) is generated by cubics, nor does it determine a complete
Markov or Graver basis.
\end{remark}
